\section*{Abstract}
This study addresses the Electric Vehicle Routing Problem (EVRP), an important extension of the classical Vehicle Routing Problem that incorporates energy and load constraints. To tackle the challenges of ensuring route feasibility under limited battery capacities while achieving high-quality solutions, three advanced metaheuristic frameworks are developed: a Genetic Algorithm (GA), a Simulated Annealing (SA) approach, and an optimized Ant Colony Optimization (ACO) method. Each framework systematically integrates perturbation-driven diversification, adaptive local search mechanisms, and problem-specific repair strategies to enhance solution quality and robustness.

Comprehensive experiments were conducted on standard EVRP benchmark sets, covering both small- and large-scale instances. The proposed algorithms were evaluated against a Greedy Search baseline under unified experimental settings, measuring total travel distance, solution feasibility, computational efficiency, and robustness over multiple independent runs.

Results demonstrate that all three metaheuristics significantly outperform the baseline, with GA consistently achieving the best performance in terms of solution quality and stability. Systematic parameter sensitivity analyses further validate the critical roles of perturbation strategies, adaptive control mechanisms, and multi-stage local search in improving search effectiveness.

Overall, this work highlights the effectiveness and scalability of perturbation-driven metaheuristic optimization for solving complex, energy-constrained routing problems and points toward promising future directions for dynamic and multi-objective EVRP scenarios.

